\documentclass[11pt]{article}

\usepackage[margin=1in]{geometry}
\usepackage{amsmath}
\usepackage{amssymb}
\usepackage{amsthm}

\theoremstyle{definition}
\newtheorem*{remark}{Remark}

\title{IE 5311-001 (D01) Principles of Optimization\\[2pt]
       \large Part 1, Problem 1: The Diet Problem (Stigler Diet)}
\author{Lidivine Kengne}
\date{\today}

\begin{document}

\maketitle

\section{Problem Statement}

Slide 15 gives the original 1945 version: for a moderately active man
weighing 154 pounds, how much of each of 77 foods should be eaten daily
so that his intake of nine nutrients is at least the recommended dietary
allowances suggested by the National Research Council in 1943, at
minimum cost?

Stigler's heuristic eliminated 62 foods and reached \$39.93. Dantzig
later solved the same instance to optimality at \$39.69 using linear
programming and the simplex method (slide 16).

The formulation below is written for arbitrary food and nutrient sets.
Substituting the full 77 foods and 9 nutrients requires changing data
only, never the model.

\section{Sets}

\[
\begin{array}{ll}
  I & \text{set of available foods, indexed by } i \\[2pt]
  J & \text{set of nutrients under control, indexed by } j
\end{array}
\]

\section{Parameters}

\[
\begin{array}{lll}
  c_i    & \text{cost per unit of food } i                    & [\text{\$}/\text{unit}] \\[2pt]
  a_{ij} & \text{amount of nutrient } j \text{ per unit of food } i & [\text{g}/\text{unit}] \\[2pt]
  \ell_j & \text{minimum required daily level of nutrient } j  & [\text{g}/\text{day}] \\[2pt]
  u_j    & \text{maximum allowed daily level of nutrient } j   & [\text{g}/\text{day}] \\[2pt]
  f_i    & \text{maximum daily servings of food } i            & [\text{unit}/\text{day}]
\end{array}
\]

All parameters are known constants, with
$c_i \ge 0$, $a_{ij} \ge 0$, and $0 \le \ell_j \le u_j$ for all
$i \in I$ and $j \in J$.

\section{Decision Variables}

\[
  x_i \;=\; \text{units of food } i \text{ consumed per day},
  \qquad x_i \in \mathbb{R}_{+}, \quad i \in I .
\]

\section{Model}

\begin{align}
  \min_{x} \quad
    & \sum_{i \in I} c_i\, x_i
      \label{eq:obj} \\[4pt]
  \text{s.t.} \quad
    & \sum_{i \in I} a_{ij}\, x_i \;\ge\; \ell_j,
      && \forall\, j \in J
      \label{eq:lower} \\[2pt]
    & \sum_{i \in I} a_{ij}\, x_i \;\le\; u_j,
      && \forall\, j \in J
      \label{eq:upper} \\[2pt]
    & 0 \;\le\; x_i \;\le\; f_i,
      && \forall\, i \in I
      \label{eq:bounds}
\end{align}

Objective \eqref{eq:obj} minimizes total daily cost. Constraints
\eqref{eq:lower} enforce the recommended minimum for every nutrient, and
\eqref{eq:upper} the corresponding maximum. Bounds \eqref{eq:bounds}
keep consumption nonnegative and cap the servings of any single food.

\section{Compact Matrix Form}

Let $x \in \mathbb{R}^{|I|}$ collect the decision variables,
$c \in \mathbb{R}^{|I|}$ the costs, $f \in \mathbb{R}^{|I|}$ the serving
caps, and $\ell, u \in \mathbb{R}^{|J|}$ the nutrient bounds. Define the
nutrient matrix $A \in \mathbb{R}^{|J| \times |I|}$ by
$A_{ji} = a_{ij}$, so that $Ax$ is the vector of daily intakes. The
model becomes
\[
  \min_{x} \; c^{\top} x
  \qquad \text{subject to} \qquad
  \ell \le A x \le u, \qquad 0 \le x \le f .
\]

\section{Classification (slide 18 taxonomy)}

\[
\begin{array}{ll}
  \text{Constrained or unconstrained}      & \text{constrained} \\
  \text{Variable type}                     & \text{continuous} \\
  \text{Variable size}                     & \text{compact, } |I| \text{ variables} \\
  \text{Objective type}                    & \text{linear} \\
  \text{Constraint type}                   & \text{linear} \\
  \text{Constraint size}                   & \text{compact, } 2|J| + 2|I| \text{ inequalities} \\
  \text{Objective size}                    & \text{single objective} \\
  \text{Parameters}                        & \text{deterministic} \\
  \text{Decision stage}                    & \text{one-stage} \\
  \text{Number of players}                 & \text{single}
\end{array}
\]

Every entry is linear with compact size, which places this problem
squarely in Part 2 of the course.

\section{Geometry (slide 19)}

Each constraint in \eqref{eq:lower} through \eqref{eq:bounds} is a
linear inequality in $x$, so each defines a half-space in
$\mathbb{R}^{|I|}$, and the boundary of each is a hyperplane. For
nutrient $j$ the equation
\[
  \sum_{i \in I} a_{ij}\, x_i = \ell_j
\]
is the hyperplane on which the minimum requirement is met exactly. The
feasible region is the intersection of $2|J| + 2|I|$ half-spaces,
therefore a polyhedron. Since $0 \le x \le f$ bounds every coordinate,
the region is a polytope: bounded, closed, and convex.

\begin{remark}[Why the serving caps matter]
Stigler's original model had no upper bounds $f_i$ and no upper bounds
$u_j$. With cost minimization as the only pressure, the optimal solution
concentrates on whichever few foods deliver nutrients most cheaply,
producing the monotonous diets the problem became known for. A solved
instance of the formulation above shows the same effect: the optimum
loads onto two foods and leaves the rest at zero. Adding $f_i$ and $u_j$
does not change the structure of the model, only its data.
\end{remark}

\begin{remark}[Reading the dual]
Each constraint in \eqref{eq:lower} carries a shadow price: the increase
in minimum daily cost caused by raising $\ell_j$ by one gram. Nutrients
whose constraints are slack at the optimum carry a shadow price of zero.
This is the dual information that slide 52 generalizes as expectation
written through dual pairing, and that Part 2 formalizes as linear
programming duality.
\end{remark}

\end{document}
